\chapter{Аналитическая часть}\label{ch:-4}

В данном разделе будут рассмотрены теоретические аспекты алгоритмов решения задачи коммивояжера.

\section{Алгоритм полного перебора}

Алгоритм полного перебора для задачи коммивояжера основывается на генерации всех возможных маршрутов, которые проходят через заданное множество городов ровно один раз. Каждый маршрут представляется как перестановка множества городов.

Процесс решения включает два этапа:
\begin{enumerate}
    \item генерация всех перестановок множества городов;
    \item для каждой перестановки вычисление общей длины маршрута на основе заданной матрицы расстояний и выбор маршрута с минимальной длиной.
\end{enumerate}

Алгоритм гарантирует нахождение оптимального решения, так как проверяются все возможные маршруты.

\section{Муравьиный алгоритм}

Муравьиный алгоритм${[2]}$ — это эвристический метод, вдохновлённый поведением колонии муравьёв при поиске оптимальных путей между источниками пищи и муравейником. Алгоритм представляет собой итеративный процесс, в котором множество муравьёв строят решения, используя информацию о феромонах и эвристическую оценку расстояний.

На каждой итерации муравей:
\begin{enumerate}
    \item сформирует 1 маршрут, который претендует на решение задачи коммивояжера на основе вероятностного правила;
    \item после завершения маршрута обновляется информация о феромонах --- на всех рёбрах испаряется часть феромона, а на рёбрах, принадлежащих маршрутам, добавляется новый феромон в зависимости от длины маршрута.
\end{enumerate}

Вероятность перехода определяется количеством феромона на ребре и эвристической полезностью.
Итерации продолжаются до выполнения заданного критерия остановки (в данном случае -- достижения максимального числа итераций).
Алгоритм не гарантирует нахождения оптимального решения.

\textbf{Модификация алгоритма по варианту:}
\begin{itemize}
    \item карта перемещения по портам через океаны (время с учётом затрат: на весь маршрут лимитирован объём средств);
    \item незамкнутый маршрут;
    \item без элитных муравьев;
    \item неориентированный граф.
\end{itemize}

\section*{Вывод}

В данном разделе были рассмотрены теоретические аспекты алгоритмов решения задачи коммивояжера.

